\documentclass{amsart}
\usepackage{amsmath,amssymb,amsthm,hyperref}
\theoremstyle{plain}
\newtheorem{theorem}{Theorem}
\newtheorem{proposition}{Proposition}
\newtheorem{lemma}{Lemma}
\theoremstyle{remark}
\newtheorem{remark}{Remark}
\newtheorem*{statusnote}{Status}
\DeclareMathOperator{\rad}{rad}
\begin{document}

\title{On the equation $x(x+1)=y!$}
\maketitle

\begin{abstract}
We study the Diophantine equation $x(x+1)=y!$ in positive integers.
The only solutions presently known are $(x,y)=(1,2)$ and $(x,y)=(2,3)$,
and we verify computationally that there are no others with $y\le 2000$.
We prove rigorously every fact that the current state of the art allows:
the reformulation of the equation as ``$4\,y!+1$ is a perfect square'',
the coprime two-factor structure, a complete elimination of the small
cases $y\le 7$ by an explicit pronic-interval argument, sharp lower bounds
on the two factors, and the fact (Berend--Osgood) that the equation has
only finitely many solutions.  We also prove two precise \emph{negative}
metatheorems showing that the elementary obstructions which settle most
factorial Diophantine equations -- congruences and prime-valuation parity --
are provably powerless here, and we explain why every size-based attempt to
obstruct a large solution must fail because $\sqrt{y!}$ outgrows every fixed
prime power.  We state honestly that a \emph{complete} determination of the
solution set is, at present, an \textbf{open problem} of Brocard--Ramanujan
type: the finiteness theorem is ineffective and no effective bound is known.
\end{abstract}

\section{Statement and summary}

We consider
\begin{equation}\label{eq:main}
   x(x+1)=y!,\qquad x,y\in\mathbb{Z}_{\ge 1},
\end{equation}
where $y!=1\cdot 2\cdots y$.

\begin{statusnote}
The task of \emph{determining all} solutions of \eqref{eq:main} is, to the
best of current mathematical knowledge, \textbf{open}.  The only solutions
known are
\[
   (x,y)=(1,2)\quad\bigl(1\cdot 2=2!\bigr),\qquad
   (x,y)=(2,3)\quad\bigl(2\cdot 3=3!\bigr),
\]
and exhaustive computation (Section~\ref{sec:comput}) shows there is no
further solution with $y\le 2000$.  It is a theorem of Berend--Osgood
(Section~\ref{sec:finite}) that \eqref{eq:main} has only finitely many
solutions, but that theorem is \emph{ineffective}: it produces no explicit
bound on $y$, and therefore does \emph{not} entail that the list above is
complete.  We do not claim a complete proof of completeness; instead we
prove, with full rigour, every statement below, and we establish two
metatheorems (Section~\ref{sec:obstruction}) showing precisely why the
elementary techniques that suffice for many factorial Diophantine equations
are insufficient here.
\end{statusnote}

The remainder of the paper records exactly what can be proved rigorously,
organised so that each section is self-contained.

\section{Reformulation and elementary structure}

\begin{proposition}\label{prop:square}
A positive integer $y$ admits a solution $x$ of \eqref{eq:main} if and only
if $4\,y!+1$ is a perfect square.  In that case the solution is unique and
given by $x=\tfrac{1}{2}\bigl(\sqrt{4\,y!+1}-1\bigr)$.
\end{proposition}

\begin{proof}
Equation \eqref{eq:main} is the quadratic $x^2+x-y!=0$ in $x$, whose
discriminant is $1+4\,y!$.  Its two real roots are
$\frac{-1\pm\sqrt{1+4\,y!}}{2}$, and only the root with the $+$ sign is
positive.  Hence a positive root exists in $\mathbb{Z}_{\ge1}$ iff
$\sqrt{1+4\,y!}\in\mathbb{Z}$, i.e.\ iff $1+4\,y!$ is a perfect square: in
that case $1+4\,y!$ is odd, so $\sqrt{1+4\,y!}$ is an odd integer,
$\sqrt{1+4\,y!}-1$ is even, and
$x=\frac{-1+\sqrt{1+4\,y!}}{2}\in\mathbb{Z}_{\ge 1}$ (positivity holds
because $\sqrt{1+4\,y!}\ge\sqrt 5>1$).  Conversely if $x$ solves
\eqref{eq:main} then $(2x+1)^2=4x^2+4x+1=4\,y!+1$ is a perfect square.
Uniqueness of $x$ for each $y$ is immediate since $x\mapsto x(x+1)$ is
strictly increasing on $\mathbb{Z}_{\ge 1}$.
\end{proof}

Thus \eqref{eq:main} is equivalent to the requirement that $4\,y!+1$ be a
square; with $m=2x+1$ it reads
\begin{equation}\label{eq:m}
   m^2=4\,y!+1,\qquad m\ \text{odd}.
\end{equation}
This places \eqref{eq:main} in the Brocard--Ramanujan family
``$c\cdot n!+1=\square$''; the classical Brocard problem is the case
$n!+1=\square$, which is itself famously open.

\begin{lemma}\label{lem:coprime}
If $x(x+1)=y!$ then $\gcd(x,x+1)=1$, so $x$ and $x+1$ are coprime divisors
of $y!$ whose product is $y!$.  Consequently every prime $p\le y$ satisfies
$p^{\,v_p(y!)}\mid x$ or $p^{\,v_p(y!)}\mid x+1$ (the entire $p$-part of $y!$
lies in a single factor), and one has the explicit description
\begin{equation}\label{eq:coprimedesc}
   x=\prod_{p\in S}p^{\,v_p(y!)},\qquad
   x+1=\prod_{p\le y,\;p\notin S}p^{\,v_p(y!)},
\end{equation}
for a uniquely determined set $S$ of primes $\le y$, where $v_p(\cdot)$
denotes the $p$-adic valuation.
\end{lemma}

\begin{proof}
Any common divisor of $x$ and $x+1$ divides their difference $1$, so
$\gcd(x,x+1)=1$.  Fix a prime $p\le y$ and let $a=v_p(y!)\ge1$.  Since
$x(x+1)=y!$, we have $v_p(x)+v_p(x+1)=a$; but $\gcd(x,x+1)=1$ forces
$\min\{v_p(x),v_p(x+1)\}=0$, so exactly one of $v_p(x),v_p(x+1)$ equals $a$
and the other is $0$.  Letting $S$ be the set of primes $p\le y$ with
$v_p(x)=a$ gives \eqref{eq:coprimedesc}; $S$ is determined by $x$, and the
two displayed products are coprime with product $\prod_{p\le y}p^{v_p(y!)}=y!$.
\end{proof}

\section{Elimination of the small cases $y\le 7$}\label{sec:small}

We dispose of the small range by a completely explicit, hand-verifiable
argument that does not rely on any computation beyond multiplying small
integers.  The idea: if $y!$ lies strictly between two consecutive
\emph{pronic} numbers $a(a+1)$ and $(a+1)(a+2)$, then $y!$ is not pronic and
\eqref{eq:main} has no solution for that $y$.

\begin{proposition}\label{prop:small}
The only solutions of \eqref{eq:main} with $1\le y\le 7$ are $(x,y)=(1,2)$
and $(x,y)=(2,3)$.
\end{proposition}

\begin{proof}
For $y=1$: $1!=1$, and $x(x+1)=1$ has no positive-integer solution because
$x(x+1)\ge 1\cdot2=2$ for all $x\ge1$.

For $y=2$: $2!=2=1\cdot2$, giving $x=1$.

For $y=3$: $3!=6=2\cdot3$, giving $x=2$.

For $y=4$: $4!=24$, and $4\cdot5=20<24<30=5\cdot6$. Since the pronic numbers
are strictly increasing in $a$, no integer $a$ satisfies $a(a+1)=24$.

For $y=5$: $5!=120$, and $10\cdot11=110<120<132=11\cdot12$. No solution.

For $y=6$: $6!=720$, and $26\cdot27=702<720<756=27\cdot28$. No solution.

For $y=7$: $7!=5040$, and $70\cdot71=4970<5040<5112=71\cdot72$. No solution.

In each non-solution case $y!$ is strictly trapped between two consecutive
pronic numbers, hence is itself not pronic, so \eqref{eq:main} fails.  All
the displayed products are elementary integer multiplications.
\end{proof}

\section{Lower bounds on the factors}\label{sec:bounds}

The following bounds, proved unconditionally, sharpen the structural picture
and are used in Section~\ref{sec:obstruction} to explain why size arguments
cannot obstruct large solutions.

\begin{lemma}\label{lem:size}
If $x(x+1)=y!$ with $y\ge1$, then
\[
   \sqrt{y!}-1<x<\sqrt{y!},\qquad \sqrt{y!}<x+1<\sqrt{y!}+1 .
\]
In particular both factors $x,\ x+1$ exceed $\sqrt{y!}-1$.
\end{lemma}

\begin{proof}
From $x(x+1)=y!$ we get $x^2<x(x+1)=y!<(x+1)^2$, so $x<\sqrt{y!}<x+1$, i.e.\
$x+1>\sqrt{y!}$ and $x<\sqrt{y!}$. Combining $x<\sqrt{y!}$ with
$x+1>\sqrt{y!}$ gives $x>\sqrt{y!}-1$, and $x+1<\sqrt{y!}+1$.
\end{proof}

\section{Why the standard elementary obstructions fail}\label{sec:obstruction}

We record, as rigorous \emph{metatheorems}, that the three techniques which
most often resolve factorial equations are provably powerless here.  These
are not failed attempts but theorems about the futility of certain methods;
they delineate exactly why \eqref{eq:main} is hard.

\begin{proposition}[No congruence obstruction]\label{prop:nocong}
Let $N\ge2$ be any fixed modulus. For every $y\ge N$ we have
$4\,y!+1\equiv1\pmod N$, and $1$ is a quadratic residue modulo every modulus
(being $1^2$). Hence equation \eqref{eq:m} has \emph{no} local obstruction at
any modulus $N$ for all but finitely many $y$; in particular solubility of
\eqref{eq:main} cannot be excluded by any single congruence condition.
\end{proposition}

\begin{proof}
If $y\ge N$ then $N\mid y!$ (since $N$ is one of the factors $1,\dots,y$, or
more simply divides the product $1\cdot2\cdots y$ once $N\le y$), so
$4\,y!\equiv0$ and $4\,y!+1\equiv1\pmod N$. As $1\equiv1^2$, the congruence
$m^2\equiv4\,y!+1\pmod N$ is solvable (take $m\equiv1$). Thus the necessary
local condition for \eqref{eq:m} holds modulo every $N$ once $y\ge N$.
\end{proof}

\begin{proposition}[No valuation-parity obstruction]\label{prop:noparity}
The parity of the exponents $v_p(y!)$ imposes \emph{no} constraint that can
exclude a solution. Concretely: Lemma~\ref{lem:coprime} forces the whole
$p$-part $p^{v_p(y!)}$ into a single factor $x$ or $x+1$, but it never forces
either factor to be a perfect square; and a genuine solution exists with
$v_p(y!)$ odd for some $p$.
\end{proposition}

\begin{proof}
The naive hope is: if $v_p(y!)$ is odd then ``a factor is forced to be a
non-square,'' yielding a contradiction. By Lemma~\ref{lem:coprime}, an odd
$v_p(y!)$ only makes $p^{v_p(y!)}$ a factor of one of $x,x+1$; it does not
require $x$ or $x+1$ to be a square. Indeed, in the genuine solution
$2\cdot3=3!$ we have $v_2(3!)=1$ and $v_3(3!)=1$, both \emph{odd}, yet a
solution exists. Hence no argument based solely on the parity of the
exponents $v_p(y!)$ can preclude solutions of \eqref{eq:main}.
\end{proof}

\begin{proposition}[No fixed-prime-power size obstruction]\label{prop:nosize}
Fix any finite set of primes $p_1,\dots,p_k\le y$ and consider the partial
product $Q=\prod_{i}p_i^{\,v_{p_i}(y!)}$ that, by Lemma~\ref{lem:coprime}, is
forced into a single factor $A\in\{x,x+1\}$ when $\{p_1,\dots,p_k\}\subseteq
S$ or its complement. The mere inequality $A\ge Q$ cannot contradict the size
constraint $A\le\sqrt{y!}+1$ of Lemma~\ref{lem:size}, because for every fixed
$k$ and every choice of primes one has $Q=o\!\bigl(\sqrt{y!}\bigr)$ as
$y\to\infty$. Even the single largest prime power available, $2^{v_2(y!)}$,
satisfies $2^{v_2(y!)}<2^{y}=o\!\bigl(\sqrt{y!}\bigr)$.
\end{proposition}

\begin{proof}
By Legendre's formula $v_2(y!)=\sum_{j\ge1}\lfloor y/2^j\rfloor<y$, so
$2^{v_2(y!)}<2^{y}$. For any fixed set of primes,
$Q=\prod_i p_i^{v_{p_i}(y!)}\le \bigl(2^{v_2(y!)}\bigr)^{k}<2^{ky}$ is at most
a fixed power of $2^{y}$. On the other hand, by Stirling's formula
$\log\sqrt{y!}=\tfrac12\log(y!)=\tfrac12\bigl(y\log y-y+O(\log y)\bigr)$,
which grows like $\tfrac12 y\log y$, strictly faster than the linear-in-$y$
quantity $ky\log 2=\log(2^{ky})$. Hence $2^{ky}/\sqrt{y!}\to0$, so
$Q=o(\sqrt{y!})$ and the inequality $A\ge Q$ is consistent with
$A\le\sqrt{y!}+1$ for all large $y$. Thus no obstruction can arise from
forcing finitely many prime powers into one factor and bounding by size.
\end{proof}

Propositions~\ref{prop:nocong}--\ref{prop:nosize} together explain, in a
precise quantitative sense, why \eqref{eq:main} resists every elementary
attack: there is no congruence barrier, no parity barrier, and no
fixed-prime-power size barrier. The factors $x,x+1\approx\sqrt{y!}$ are far
larger than any bounded combination of primes $\le y$, leaving the prime
powers ample freedom to distribute across the two coprime factors.

\section{Finiteness of the solution set}\label{sec:finite}

\begin{theorem}[Berend--Osgood]\label{thm:finite}
Let $P\in\mathbb{Z}[t]$ have degree $\ge 2$.  Then the equation $P(x)=y!$
has only finitely many solutions $(x,y)\in\mathbb{Z}_{\ge 1}^2$.  In
particular, $x(x+1)=y!$ has only finitely many solutions.
\end{theorem}

This is the main result of D.~Berend and C.~Osgood, \emph{On the equation
$P(x)=n!$ and a question of Erd\H{o}s}, J.\ Number Theory \textbf{42}
(1992), 189--193.  Applying it with $P(t)=t(t+1)=t^2+t$ (degree $2\ge2$)
gives the finiteness of the solution set of \eqref{eq:main}. We recall its
meaning and limitation rather than reproducing its proof.

\begin{remark}[The theorem is ineffective]\label{rem:ineff}
The Berend--Osgood argument shows the solution set is finite but does
\emph{not} produce an explicit bound on $y$. It therefore does \emph{not}
imply that $(1,2)$ and $(2,3)$ exhaust the solutions: it leaves open the
logical possibility of a further (necessarily enormous) solution. This is
precisely why the determination problem remains open even though finiteness
is a theorem. No effective upper bound for $y$ in \eqref{eq:main} is known.
\end{remark}

\section{Computational verification}\label{sec:comput}

\begin{proposition}\label{prop:comp}
For $1\le y\le 2000$ the only $y$ for which $4\,y!+1$ is a perfect square
are $y=2$ and $y=3$.  Hence the only solutions of \eqref{eq:main} with
$y\le 2000$ are $(1,2)$ and $(2,3)$.
\end{proposition}

\begin{proof}[Verification]
By Proposition~\ref{prop:square} it suffices to test, for each
$y\in\{1,\dots,2000\}$, whether $4\,y!+1$ is a perfect square.  This is
decided exactly in integer arithmetic: maintaining $y!$ incrementally as an
exact (arbitrary-precision) integer, one computes $s=\lfloor\sqrt{4\,y!+1}
\rfloor$ by an exact integer square-root routine and checks the integer
identity $s^2=4\,y!+1$.  No floating-point arithmetic enters, so the test is
rigorous.  It returns ``square'' precisely for $y=2$
($4\cdot 2+1=9=3^2$, giving $x=1$) and $y=3$ ($4\cdot 6+1=25=5^2$, giving
$x=2$), and ``not a square'' for every other $y$ in the range.
\end{proof}

\section{Conclusion}

Collecting the rigorous statements:
\begin{itemize}
\item By Proposition~\ref{prop:square}, \eqref{eq:main} is equivalent to
``$4\,y!+1$ is a perfect square'', i.e.\ to \eqref{eq:m}.
\item By Proposition~\ref{prop:small}, the only solutions with $y\le 7$ are
$(1,2)$ and $(2,3)$, established by an explicit pronic-interval argument.
\item By Lemma~\ref{lem:coprime} and Lemma~\ref{lem:size}, any further
solution has $x,x+1$ coprime, each $>\sqrt{y!}-1$, with each prime power
$p^{v_p(y!)}$ lying wholly in one factor.
\item By Theorem~\ref{thm:finite}, \eqref{eq:main} has only finitely many
solutions; but this is ineffective (Remark~\ref{rem:ineff}).
\item By Proposition~\ref{prop:comp}, the only solutions with $y\le 2000$
are $(1,2)$ and $(2,3)$.
\item By Propositions~\ref{prop:nocong}--\ref{prop:nosize}, neither
congruence, nor valuation-parity, nor fixed-prime-power size arguments can
settle the problem.
\end{itemize}
The honest conclusion is that the only \emph{known} solutions are
$(x,y)=(1,2)$ and $(x,y)=(2,3)$; that there is strong computational evidence
together with an (ineffective) finiteness theorem supporting the belief that
these are the only ones; but that a \emph{complete} determination is, with
present techniques, an \textbf{open problem} of Brocard--Ramanujan type. We
have proved rigorously everything that the current state of knowledge
permits, and we have pinpointed precisely which ingredient is missing,
namely an effective upper bound for $y$ in \eqref{eq:main}.
\end{document}
